\emph{(i).} Clause (v) of Theorem \claimref{synthesis.gravitational-pair-production-from-the-operator-source-at-first-order} gives $\alpha_{j} = \tfrac{\pi^{3}}{17280}\bigl[\bigl(c_{j}(N_{s} + 12N_{v}) + 5\tau_{j}\tau_{B}\bigr)/225 + 49c_{j}N_{w}/4800\bigr]$ with $c_{j} = 1, 3, 12$ and $\tau_{j} = 0$ unless $j$ is a scalar. The expression is linear in $(c_{j}, \tau_{j})$, so the sum over $D$ replaces $c_{j}$ by $C_{D}$ and $\tau_{j}$ by $\tau_{D}$. The term $5\tau_{j}\tau_{B}$ comes from the weights $(\pi/180)(c_{i}c_{j} + 5\tau_{i}\tau_{j})$ of clause (iv) summed over the bath's scalars; that clause states the weight for species with their own couplings, and the script verifies it directly for scalars to scalars with two different couplings: the bracket of the two stored pair sources is $\tfrac{s^{2}}{4}(\cos^{2}\theta - \tfrac13) - \tfrac{s^{2}}{6}(1 - 6\xi_{1})(1 - 6\xi_{2})$ at five rational angles, and the integral of its square over the sphere, with the factor $\tfrac12$ for each pair of identical quanta, is $(\pi/180)(1 + 5\tau_{1}\tau_{2})$ at seven pairs of couplings. The three consequences are read off: $\tau_{D} = 0$ for a content without scalars; the product vanishes when either factor does; and $\tau_{D}\tau_{B} \ge 0$ with each $\tau$ increasing in $|1 - 6\xi|$.

\emph{(ii).} With $C_{D} = 4 + 3 \cdot 45 + 12 \cdot 12 = 283$ and the bath $(4, 45, 12)$, $c/\pi^{3} = \tfrac{1}{17280}\bigl[\tfrac{283 \cdot 148}{225} + \tfrac{49 \cdot 283 \cdot 45}{4800}\bigr] + \tfrac{5\tau_{D}\tau_{B}}{225 \cdot 17280}$, and $5/(225 \cdot 17280) = 1/777600$; at $\tau_{D} = \tau_{B} = 4$ it is the stored $1519247/82944000$. The stored readings of the couplings are those of Theorem \claimref{synthesis.no-single-curvature-coupling-serves-the-sakharov-zero-and-the-kasner-tilt}: $0$ for the doublets, $\tfrac16$ for the attachment's helicity-zero module. At $\tau_{D}\tau_{B} = 0$ the value is lower by $16/777600$, which is $5120/4557741$ of the stored one; in numbers, with $103993/33102 < \pi < 355/113$, $0.5672 < c < 0.5673$ against $0.5679 < c < 0.5680$. The threshold $c^{*}$ is recomputed from the inputs of Proposition \claimref{synthesis.seeding-remainder-excludes-perturbative-pati-salam} ($\lfloor 10^{4}c^{*} \rfloor = 134949$); the margin $1.56$ at four contents and the count $\lfloor c^{*}/c \rfloor = 23$ hold at both values of $c$. The window is the stored one of Theorem \claimref{synthesis.tree-level-curvature-couplings-stay-blue-at-the-isotropic-kasner-point}; $\tau = (1 - 6\xi)^{2}$ at its ends is $66.186$ and $68.699$, and $\tau_{D}\tau_{B} = 16\tau^{2}$. The stated values of $c$, the ratios, the position of the crossing (below $c^{*}$ at $\xi = 1.5239$ with $\pi$ at its upper end, above it at $1.5241$ with $\pi$ at its lower end), the central $c^{*}$ (remainder $4.364$, $8.7 \times 10^{-11}$ baryons per unit entropy, $\kappa_{s} = 145.26$, $v_{R} = 4.2 \times 10^{10}$ GeV) and the three cases of one further real scalar (bath $(5, 45, 12)$ with $\tau_{B} = 4 + \tau$; the same with the copy enlarged to $C_{D} = 284$, $\tau_{D} = 4 + \tau$; and $C_{D} = 1$, $\tau_{D} = \tau$ from the Standard Model's bath) are evaluations of the formula of (i), made by the script in exact arithmetic with every margin compared in fourth powers.

\emph{(iii).} The content between $M_{I}$ and $M_{U}$ is taken from Theorem \claimref{synthesis.series-content-induces-a-newton-coupling-only-through-an-unfixed-cut-off}(ii), which quotes Theorem \claimref{standard-model.chain}: twenty-one vectors, forty-eight Weyl fields, and scalars whose number is fixed by the stored values $-36 + N_{s} = 28, 32, 152$ of the induced coefficient at minimal coupling, that is $N_{s} = 4 + 60$, $8 + 60$, $8 + 60 + 120$. The clause's largest $M_{3}$ is $4\pi v_{R} = 5.34 \times 10^{11}$ GeV, and $1.9 \times 10^{10} < 5.34 \times 10^{11} < 6 \times 10^{14} < 1.1 \times 10^{16}$ GeV (the last being the least of the stored unification scales), so on the chain a bath at the clause's temperatures above $M_{I}$, or at the ceiling, has this content; below $M_{I}$ the bath is the Standard Model's and the temperature is lower still, so the clause's comparison is weakest above $M_{I}$. The remaining statements are evaluations of the formula of (i) with $(N_{s}, N_{w}, N_{v}) = (N_{s}, 48, 21)$: a vector has $C_{D} = 12$, $\tau_{D} = 0$; a Standard-Model content has $C_{D} = 283$, $\tau_{D} = 4$ with $\tau_{B} = N_{s}$; a copy of the chain's content has $C_{D} = N_{s} + 3 \cdot 48 + 12 \cdot 21$ and $\tau_{D} = \tau_{B} = N_{s}\tau$. For $N_{s} = 188$ and one coupling, $c$ is an increasing function of $\tau^{2}$, so $c < c^{*}$ on an interval of $|1 - 6\xi|$ about zero; the script finds $c < c^{*}$ at $\xi = 0.022$ and $0.311$ and $c > c^{*}$ at $0.021$ and $0.313$, at the ends of the bracket of $\pi$ that weaken each statement, and in the same way for $N_{s} = 64$ and $68$ that $c < c^{*}$ at $\xi = 0.46$ and $c > c^{*}$ at $0.48$.

\emph{The Reading} is argued in the statement from the texts of the stored results; it is not used above.
